\section{Guarantees}\label{sec:guarantees}

This section states which of CONSENT's guarantees transfer unchanged,
which require restatement under persistence and the committed FSL, and
what is newly provable about the commitment itself. Throughout, $J_d$
denotes the day-$d$ building cost under the charging policy of
Section~\ref{sec:charging}, inclusive of the FSL terms \eqref{eq:dr}
when an event is active, and $\theta_v$, $\bm{\theta}_{\setminus v}$
are as in \eqref{eq:uflex}.

\subsection{Standing assumptions}

\begin{assumption}[Fixed assignment]\label{as:assign}
The charger assignment $\assign$ is fixed at arrival for the entire
session (Section~\ref{sec:specs}); a report does not alter the
assignment of any other user.
\end{assumption}

\begin{assumption}[Oblivious users]\label{as:oblivious}
Users best-respond to the posted menu of the current session, taking
the forecast model $\psihat$ and the choices of others as exogenous;
they do not act to manipulate the system's future state or its learned
forecasts. This is the mechanism-side counterpart of the physical
exogeneity already assumed of arrivals and off-site usage
(Assumption~\ref{as:exo}).
\end{assumption}

\begin{assumption}[Controller-set banking]\label{as:bankctrl}
The banked amount and its timing are controller decisions; the user's
only banking action is the binary consent $\accept{v}{k}$
(Section~\ref{sec:negotiation}).
\end{assumption}

\begin{assumption}[Effective floor]\label{as:floor}
$\Ksoc$ is large enough that the negotiated floor \eqref{eq:floor}
binds as a hard service constraint at the negotiated departure.
\end{assumption}

\subsection{Incentive compatibility under persistence and DR}

The single-shot strategy-proofness of CONSENT does not survive
persistence unqualified: a returning user's valuation of departure SoC
includes the shadow value of charge carried to the next session, a
quantity the single-day mechanism neither elicits nor prices. The
banking rate \eqref{eq:pbank} restores the guarantee up to a forecast
error, and the DR layer adds no new deviation because no user-facing
quantity depends on a user's report of event-related information.

\begin{lemma}[Objective-agnostic monotonicity]\label{lem:setincl}
Let $X(\theta_v)$ be the feasible action set of the day-$d$ charging
problem given $v$'s report, and let $G$ be any cost functional. If
$X(\bar\theta_v) \subseteq X(\theta_v)$, then
$\min_{x \in X(\bar\theta_v)} G(x) \ge \min_{x \in X(\theta_v)} G(x)$.
\end{lemma}

\begin{proof}
A minimum over a subset is no smaller than the minimum over the
superset.
\end{proof}

The lemma is what makes the incentive argument survive every change of
objective in this paper, including the FSL penalty: the proof of
CONSENT's Theorem~1 uses only feasible-set inclusion, never the form of
the cost.

\begin{theorem}[Periodic ex-post IC up to forecast error]\label{thm:ic}
Under Assumptions~\ref{as:assign}--\ref{as:floor}, on every day $d$,
for every history and every realization of
$\bm{\theta}_{\setminus v}$, a user cannot gain more than
$\epsIC{v}{d}$ by misreporting an earlier departure time or a higher
required SoC:
\[
\Uflex(\theta_v) \;\ge\; \Uflex(\bar\theta_v) \;-\; \epsIC{v}{d},
\qquad
\epsIC{v}{d} \;=\; \bigl(c^{\mathrm{next}}_{v} -
\cmarg{v}\bigr)^{+}\, \banked{v}{k},
\]
where $c^{\mathrm{next}}_{v}$ is the realized next-session marginal
cost. If forecasts are exact, the mechanism is exactly periodic
ex-post IC; if additionally $v$ has a single session, the statement
reduces to CONSENT's Theorem~1.
\end{theorem}

\begin{proof}
\emph{Step 1 (report channel).} Reporting
$\bar{t}_{\mathrm{req}} < \treq$ requires any feasible schedule to set
$\rate{t}{v} = 0$ after $\bar{t}_{\mathrm{req}}$ and to satisfy the
floor by $\bar{t}_{\mathrm{req}}$; every such schedule is feasible
under the true report (Assumption~\ref{as:assign} rules out
reassignment effects). Reporting $\bar{E}_{\mathrm{req}} > \Ereq$
intersects the feasible set with a smaller terminal set
(Assumption~\ref{as:floor}). In both cases
$X(\bar\theta_v) \subseteq X(\theta_v)$. The day-$d$ cost, with or
without the event terms \eqref{eq:dr}, is a fixed functional of the
action because $\FSL_d$ was committed before any arrival. By
Lemma~\ref{lem:setincl}, for every realization and every sampled
scenario,
$J_d(\bm{\theta}_{\setminus v}, \bar\theta_v) \ge
 J_d(\bm{\theta}_{\setminus v}, \theta_v)$.
Since the leave-out term
$J_d(\bm{\theta}_{\setminus v})$ in \eqref{eq:uflex} is common,
$\Uflex(\theta_v) \ge \Uflex(\bar\theta_v)$ pointwise.

\emph{Step 2 (banking channel).} By
Assumption~\ref{as:bankctrl} the only remaining action is
$\accept{v}{k}$. The gain from consenting with departing excess
$\banked{v}{k}$ is
\[
c^{\mathrm{next}}_{v}\,\banked{v}{k} - \pbank{v}{k}\,\banked{v}{k}
\;\le\; \bigl(c^{\mathrm{next}}_{v} - \max(\bar{w}^{e}_{k},
\cmarg{v})\bigr)\,\banked{v}{k}
\;\le\; \epsIC{v}{d},
\]
using \eqref{eq:pbank} with $\muDR = 0$ on ordinary days (on
pre-event days the subsidy $\muDR\rhohat_v$ is a deliberate transfer
for buffer the building recovers in the event window, and the same
bound applies to the residual). The $\Ubuf$ payment cannot be inflated
by any report: it is settled ex post by replay \eqref{eq:ubuf} on
realized quantities only.

\emph{Step 3 (no DR channel).} The user's action set contains no
event-indexed report and no quantity on which $\FSL_d$, $\winc$, or
$\Kfsl$ conditions; the DR layer alters only the coefficients of the
day-$d$ cost, already covered by Step~1. Under
Assumption~\ref{as:oblivious} no strategy conditions on $\psihat$.
Combining Steps~1--2 gives the claim.
\end{proof}

\begin{remark}
The guarantee is deliberately weaker than dominant-strategy: once
private information evolves across sessions, periodic ex-post IC up to
a model-error term is the strongest notion attainable without full
dynamic-mechanism machinery, and $\epsIC{v}{d}$ is empirically
reportable from the forecast diagnostics of Section~6.
\end{remark}

\subsection{Budget feasibility and participation}

\begin{theorem}[Per-day budget feasibility of flexibility
payments]\label{thm:bf}
With $\ucost{k}{v}$ as in \eqref{eq:usercost} and
$\alpha \in [0,1]$: (i) whenever $\Uflex(\theta_v) \ge 0$, serving $v$
at the negotiated terms costs the building no more than excluding $v$,
on event and non-event days alike; (ii) in aggregate,
$\mathbb{E}\bigl[\sum_{v} \alpha\,\Uflex(\theta_v)\bigr] \le
\mathbb{E}\bigl[J_d(\emptyset) - J_d(\bm{\theta})\bigr]$, i.e., total
flexibility payments never exceed total expected savings.
\end{theorem}

\begin{proof}
(i) The building's day-$d$ ledger with $v$ satisfies
$J_d(\bm{\theta}_{\setminus v}, \theta_v) + \alpha\,\Uflex(\theta_v)
\le J_d(\bm{\theta}_{\setminus v}, \theta_v) + \Uflex(\theta_v) =
J_d(\bm{\theta}_{\setminus v})$ by \eqref{eq:uflex}. The capacity
payment $\Cinc{d}$ is fixed by $\FSL_d$ and cancels between the two
sides; $v$'s event-day contribution enters $\Uflex$ only through
penalty reduction, a genuine saving, so the inequality holds per day.
(ii) $\Uflex$ is computed sequentially at each arrival, conditional on
realized predecessors and in expectation over successors; by the tower
property the sum telescopes,
$\mathbb{E}[\sum_v \Uflex] = \mathbb{E}[J_d(\emptyset) -
J_d(\bm{\theta})]$, and $\alpha \le 1$ gives the bound. The
telescoping is what rescues the aggregate from the usual objection
that marginal values need not sum to the total.
\end{proof}

\begin{lemma}[Budget feasibility of replay payments]\label{lem:replaybf}
Let $\Ubuf(v,k)$ be the replay marginal \eqref{eq:ubuf} and let the
paid amount be the capped marginal of
Section~\ref{sec:negotiation}. Then
$\sum_{v} \text{paid}(v) \le J^{\mathrm{re}}_d(\emptyset) -
J^{\mathrm{re}}_d(\buff{d})$: buffer payments never exceed the
realized saving the buffer produced. If $J^{\mathrm{re}}_d$ is
submodular in the buffer set, the cap never binds and full marginals
are paid.
\end{lemma}

\begin{proof}
The proportional-share cap allocates exactly the realized total saving,
so the capped sum cannot exceed it. Under submodularity, marginal
contributions of a set of resources under-sum the total saving, so
each plain marginal is below its proportional counterpart's aggregate
constraint and the $\min$ selects the marginal. Substitutability of
buffers shaving a common peak is the submodular case; complementary
buffers are exactly where the cap binds.
\end{proof}

Voluntary participation transfers unchanged: the reject-all option and
the binary banking consent bound every user's exposure by the external
rate $\wext$, per session, and event days only enlarge $\Uflex$ at
fixed $\alpha$, weakly improving every offered option.

\subsection{The commitment: optimality anchor and guardrail}

\begin{proposition}[Newsvendor characterization of the risk-neutral
commitment]\label{prop:newsvendor}
Consider the LP relaxation of the day-$d$ delivery problem and let
$Q(\FSL, \omega)$ be its optimal cost in scenario $\omega$, inclusive
of \eqref{eq:dr}. Then $Q(\cdot, \omega)$ is convex, and the
risk-neutral commitment $\FSL^{\star} = \arg\min_{\FSL}
\bigl\{-\winc(\Phist - \FSL) + \mathbb{E}[Q(\FSL,\omega)]\bigr\}$
satisfies, in subdifferential form,
\[
\winc \;\in\; \Kfsl\,\dt\; \mathbb{E}_{\omega}\bigl[\,
\Nexc(\FSL^{\star}, \omega)\,\bigr],
\]
where $\Nexc(\FSL,\omega) = |\{t \in \evwin{d} :
\Ptot{t}(x^{\star}(\FSL,\omega)) > \FSL\}|$ counts exceedance steps
under the optimized recourse. The certain per-kW capacity rate equals
the expected marginal penalty over exceedance steps; the cost of
procuring user flexibility shapes the condition only through the
distribution of $\Nexc$, not additively, because the program imposes a
soft penalty rather than a hard cap.
\end{proposition}

\begin{proof}
$\Ptot{t}(x)$ is affine in the rates, so $(\Ptot{t}(x) - \FSL)^{+}$ is
jointly convex in $(x, \FSL)$; the remaining cost is convex in $x$
under the LP relaxation; partial minimization over $x$ preserves joint
convexity, so $Q(\cdot,\omega)$ is convex, and so is the bid objective
after adding the affine revenue term. By Danskin's theorem,
$\partial_{\FSL} Q(\FSL,\omega) \ni -\Kfsl\,\dt\,\Nexc(\FSL,\omega)$
almost everywhere (at ties, the convex hull over optimizers).
Stationarity of the bid objective yields the displayed condition;
convexity makes it globally sufficient. Sanity: with
$\winc/(\Kfsl\,\dt)$ exceedance steps tolerated at the optimum, an
interior solution exists whenever this ratio is below
$|\evwin{d}|$, which holds at the program rates and a $4$--$5$\,h
window.
\end{proof}

\begin{remark}[MILP caveat]\label{rem:milp}
The value function of the mixed-integer delivery problem is not convex
in $\FSL$ in general; Proposition~\ref{prop:newsvendor} is stated for
the LP relaxation, and convexity of the integer value over the
operating range of $\FSL$ is verified numerically in Section~7. The
depth $\lambda^{\star}$ implied by $\FSL^{\star}$ through
\eqref{eq:pledgedepth} anchors the deployed pair
$(\lambda, \varepsilon)$: operating at $\lambda < \lambda^{\star}$ is
priced risk aversion.
\end{remark}

\begin{proposition}[Demand-charge guardrail]\label{prop:guardrail}
Let $R_d$ be the total EV energy that must be delivered on event day
$d$, let
$E_{\mathrm{win}}(\FSL) = \dt \sum_{t \in \evwin{d}}
(\FSL - \bload{t})^{+}$
be the maximum EV energy deliverable inside the window at zero
penalty, and let
$C_{\mathrm{out}} = \dt \sum_{t \notin \evwin{d}}
\min\bigl(r^{\mathrm{cap}}_t,\ (\pmaxpast - \bload{t})^{+}\bigr)$
be the energy deliverable outside the window without raising the
billing-cycle peak, where $r^{\mathrm{cap}}_t$ aggregates assigned
charger limits. If $R_d - E_{\mathrm{win}}(\FSL) \le
C_{\mathrm{out}}$, a zero-penalty schedule exists that does not raise
the monthly peak. Since $E_{\mathrm{win}}$ is non-decreasing,
$\Fsafe = \min\{\FSL \ge 0 : R_d - E_{\mathrm{win}}(\FSL) \le
C_{\mathrm{out}}\}$
is well defined, and any commitment $\FSL_d \ge \Fsafe$ avoids the
conflict between the commitment and the demand charge
deterministically.
\end{proposition}

\begin{proof}
A zero-penalty schedule delivers at most $E_{\mathrm{win}}(\FSL)$
inside the window, so at least $R_d - E_{\mathrm{win}}(\FSL)$ must
move outside; $C_{\mathrm{out}}$ is exactly the out-of-window capacity
available under the ceiling $\pmaxpast$ and the charger limits, and
feasibility of the split follows by filling out-of-window steps
greedily up to their individual caps (supply at least demand with
per-step capacities). Monotonicity of $E_{\mathrm{win}}$ is
term-by-term.
\end{proof}

Committing below $\Fsafe$ is not forbidden; it is the regime in which
the expected capacity gain of Proposition~\ref{prop:newsvendor} must
be weighed against a certain demand-charge increment, a one-line
comparison once both are computed.

\subsection{Conditional temporal decomposition}

\begin{theorem}[Decomposition within the target-coupled policy
class]\label{thm:decomp}
Let $\sigma_d = (\pmaxpast, \{\Earr(v,k_d)\}_{v})$ and let
$\Pi_{\mathrm{CE}}$ be the class of policies consisting of per-day
charging and negotiation policies together with the terminal targets
$\{\etarget{v}\}$ of \eqref{eq:bshort}. Then: (a) $\sigma_d$ is a
sufficient statistic for the past; (b) within $\Pi_{\mathrm{CE}}$, the
billing-cycle problem separates into per-day problems coupled only
through $\sigma_d$, the demand charge decomposing as a maximum over
days exactly as in the non-persistent case; (c) the optimality gap of
$\Pi_{\mathrm{CE}}$ relative to unrestricted policies is nonnegative
and vanishes when $\psihat$ is exact, and it is the value of closed-loop
revision of the targets otherwise.
\end{theorem}

\begin{proof}
(a) The exogenous processes are independent of control by assumption
(Section~\ref{sec:specs}); the only controlled quantities that persist
across a day boundary are the running peak and the persistent SoC,
both carried in $\sigma_d$. (b) The billing-cycle maximum over
disjoint daily step sets equals the maximum of daily maxima, carried
by $\pmaxpast$. Within $\Pi_{\mathrm{CE}}$, day $d{-}1$ influences day
$d$ only through $\Edep(v, k_{d-1})$, which the terminal reward
controls; conditional on realized $\sigma_d$, the day-$d$ problem
shares no decision variable with any other day. (c) With $\psihat$
exact, the commitment layer's recourse computes the true event-day
problem; setting $\etarget{v}$ to its optimal pre-event departure SoC
and $\Kbank{v}$ large reproduces the unrestricted optimum by
concatenation, so the gap vanishes. In general the gap is the
certainty-equivalence loss of freezing the targets a day ahead, which
Section~7 reports numerically.
\end{proof}

The theorem scopes precisely what the architecture assumes: the
cross-day coupling is summarized by one scalar target per user, and
the price of that summary is a measurable certainty-equivalence gap,
not a hidden approximation.
